\documentclass[10pt]{article}

\usepackage[margin=0.6in]{geometry}
\usepackage[T1]{fontenc}
\usepackage{lmodern}
\usepackage{microtype}
\usepackage{booktabs}
\usepackage{tabularx}
\usepackage{enumitem}

\pagestyle{empty}
\setlength{\parindent}{0pt}
\setlength{\parskip}{2pt}
\setlist{leftmargin=1.1em, itemsep=0.5pt, parsep=0pt, topsep=1pt}
\newcolumntype{Y}{>{\raggedright\arraybackslash}X}
\newcommand{\head}[1]{\par\vspace{4pt}{\large\bfseries #1}\par\vspace{1pt}}
\newcommand{\person}[2]{\par\vspace{2pt}\textbf{#1} \textit{(planned: #2)}}

\begin{document}

\begin{center}
  {\large\bfseries ECE473 Milestone I Update: Replicating USDA Rainfall
  Insurance on Kalshi}\\[1pt]
  Maya, Evan, and Jack \quad\textbullet\quad October 29, 2026
\end{center}
\vspace{-8pt}

\textbf{Summary.} All nine Milestone I tasks in our one-pager are done. Today's
Kalshi monthly rain contracts capture 43\% of the variation in USDA Pasture,
Rangeland, Forage (PRF) insurance payouts at our seven study counties, at about
1.8 times PRF's cost to a rancher per dollar of expected payout.

% ---------------------------------------------------------------------------
\head{1. Planned vs.\ actual progress}

\person{Jack}{farm exposures; weather data and the CPC-to-NCEI change;
geographic basis risk}
\begin{itemize}
  \item \textbf{W1:} built 11 hay and grazing farms in 7 PRF counties, each 29
    to 55 miles from a Kalshi station.
  \item \textbf{W2:} rebuilt the NOAA CPC index PRF paid on through 2026,
    matching 854 published finals within 0.002 points. \emph{Challenge:}
    USDA's tool now shows only a recomputed NCEI history, 12.6 points off on
    average.
  \item \textbf{W3:} station rain correlates only 0.74 to 0.89 with the farm
    grid. Only 7.4\% of PRF grids are within 50 miles of a station.
\end{itemize}

\person{Maya}{PRF rules and payout math; Kalshi markets, cost, and trading;
replication quality}
\begin{itemize}
  \item \textbf{W1:} built a PRF calculator that reproduces the USDA handbook
    example and 73 published payment factors. \emph{Challenge:} the 2027 rules
    changed allocation caps mid-project.
  \item \textbf{W2:} documented rain ladders for about 35 cities, the fee
    schedule, and a longshot bias. \emph{Challenge:} five of our seven cities first
    listed on 10/25/2026, so only Dallas and Denver have price history.
  \item \textbf{W3:} hedge effectiveness falls from 1.00 to 0.73 (station
    gauge), 0.54 (monthly), and 0.43 (whole-inch strikes).
\end{itemize}

\person{Evan}{farmer scenarios; settlement proof on each side; PRF vs.\ Kalshi
payouts}. We also added rain-yield effects by crop and how farmers buy
insurance.
\begin{itemize}
  \item \textbf{W1:} modeled premiums, payouts, and uncovered losses for 11
    farms. The 90th-percentile uncovered loss is \$1,800 to \$24,700 a year.
  \item \textbf{W2:} PRF pays on gridded data, Kalshi on one NWS airport
    report. Our station data reproduces all 922 settled Kalshi outcomes.
  \item \textbf{W3:} Kalshi matched only 70\% of PRF dollars in the same year
    (2015--2025) and costs \$1.09 per \$1 of expected payout, versus \$0.61 for
    subsidized PRF. \emph{Challenge:} order books fill only 10 to 17\% of a
    Dallas hedge.
\end{itemize}

% ---------------------------------------------------------------------------
\head{2. Remaining work for Milestone II (due 11/12)}

{\small
\begin{tabularx}{\textwidth}{@{}lYY@{}}
\toprule
Owner & Week 4: designing better options & Week 5: building the app \\
\midrule
Jack & Optimizer for the lowest-error, lowest-cost portfolio, with fees,
  capital lock-up, and depth limits & Live-price backend that quotes a
  farmer's custom plan \\
Maya & Split the remaining error into binning, geography, settlement,
  liquidity, and capital & Farmer screen: location to PRF grid and nearest
  station, plus months to cover \\
Evan & Find exposures current contracts replicate poorly and propose contracts
  that fix them & Test the app on the 11 farms against the PRF benchmark \\
\bottomrule
\end{tabularx}
}
\par\vspace{1pt}
All three of us will also fact-check six research notes and re-pull order books
after 11/1.

% ---------------------------------------------------------------------------
\head{3. Refinements to our plan}
\begin{itemize}
  \item \textbf{Scope:} complement PRF instead of replacing it. The subsidy is
    the whole cost gap, so we will target coverage above PRF's 90\% cap,
    uncovered intervals, and excess-rain risk.
  \item \textbf{Approach:} the optimizer will cap each contract at order-book
    depth, which binds before price. Backtests use the as-paid CPC index, and
    forward results use NCEI.
  \item \textbf{App:} the farmer screen will show station distance and expected
    hedge effectiveness, so basis risk is visible before buying.
  \item \textbf{Timeline:} the fact-check and order-book re-pull move to the
    start of Week 4, since our depth snapshot came during Kalshi's liquidity
    rewards. \textbf{Responsibilities:} no changes.
\end{itemize}

\end{document}
